\section{Lineární kódy}
\subsection{Motivace -- bezpečnostní kódy}
Bezpečnostní kódy odhalují, zda při přenosu nedošlo k chybě, případně chyby opravují. K informačním znakům musíme přidat kontrolní znaky. Čím více chyb chceme odhalit, tím více kontrolních znaků musíme přidat. Chceme sestrojit kódy, které budou odhalovat daný počet chyb a přitom budou mít rozumně nízkou redundanci. Navíc chceme, aby tyto kódy měly jednoduchý algoritmus opravování chyb.

Pokud na konečné kódové abecedě $T$ budeme umět sčítat a násobit tak, aby $T$ bylo těleso, pak množina $T^n$ všech slov délky $n$ bude lineární prostor a budeme moci použít aparát lineární algebry.

Lineární prostory nad obecným tělesem se \enquote{chovají podobně} jako lineární prostory nad tělesem reálných čísel.
